% Stage II deliverable: Raport z testów penetracyjnych (MediSlot / NovaMed)
% Build: pdflatex (uruchom dwukrotnie dla spisu treści i odnośników).
\documentclass[11pt,a4paper]{article}

% --- Kodowanie i język ---
\usepackage[utf8]{inputenc}
\usepackage[T1]{fontenc}
\usepackage{lmodern}
\usepackage[polish]{babel}

% --- Układ ---
\usepackage[a4paper,margin=2.5cm]{geometry}
\usepackage{setspace}
\usepackage{parskip}

% --- Tabele ---
\usepackage{booktabs}
\usepackage{tabularx}
\usepackage{longtable}
\usepackage{array}
\usepackage{ragged2e}

% --- Grafika / diagramy ---
\usepackage{xcolor}
\usepackage{graphicx}
\usepackage{tikz}
\usetikzlibrary{positioning,fit,backgrounds,arrows.meta,shapes.geometric}

% --- Kod / listingi ---
\usepackage{listings}
\lstset{
  basicstyle=\ttfamily\footnotesize,
  breaklines=true,
  columns=fullflexible,
  frame=single,
  framesep=4pt,
  backgroundcolor=\color{gray!6},
  keepspaces=true,
}

% --- Różne ---
\usepackage{enumitem}
\usepackage{fancyhdr}
\usepackage[hidelinks]{hyperref}

% --- Kolumny dla tabularx ---
\newcolumntype{L}{>{\RaggedRight\arraybackslash}X}
\newcolumntype{C}{>{\Centering\arraybackslash}X}

% --- Paleta istotności ---
\definecolor{crit}{HTML}{B00020}
\definecolor{high}{HTML}{C8651B}
\definecolor{med}{HTML}{B38600}
\definecolor{low}{HTML}{3A7D44}
\newcommand{\sev}[2]{\textbf{\textcolor{#1}{#2}}}

% --- Placeholder na zrzut ekranu ---
% Użycie: \zrzut{opis treści}{IDENTYFIKATOR}
% Po zebraniu dowodu zastąp wnętrze \fbox przez: \includegraphics[width=0.92\textwidth]{sciezka}
\newcommand{\zrzut}[2]{%
  \begin{center}%
  \fbox{\begin{minipage}[c][30mm][c]{0.92\textwidth}%
    \centering\itshape [ MIEJSCE NA ZRZUT EKRANU --- #1 ]\\[2pt]%
    \normalfont\footnotesize Identyfikator: \texttt{#2}. Sposób pozyskania: patrz \emph{Instrukcja pozyskania dowodów} (osobny dokument), poz. \texttt{#2}.%
  \end{minipage}}%
  \end{center}}

% --- Nagłówek / stopka ---
\pagestyle{fancy}
\fancyhf{}
\fancyhead[L]{MediSlot --- NovaMed Sp. z o.o.}
\fancyhead[R]{Raport z testów --- Etap II}
\fancyfoot[C]{\thepage}
\renewcommand{\headrulewidth}{0.4pt}

\setlength{\tabcolsep}{5pt}
\renewcommand{\arraystretch}{1.25}

\title{\textbf{Test penetracyjny środowiska MediSlot}\\[0.3em]
\large Raport z przebiegu testów --- Etap II\\[0.2em]
\normalsize NovaMed Sp. z o.o.}
\author{Zespół \\ Kamil Szkarłat \textbullet{} Bartłomiej Ząbek \textbullet{} Jakub Szkaradek \textbullet{} Radosław Chruściński}
\date{26.11.2026 r.}

\begin{document}

\maketitle
\begin{center}
\emph{Dokument wewnętrzny, laboratorium dydaktyczne. Dane wyłącznie syntetyczne. Środowisko uruchamiane lokalnie (Docker Compose), nie wystawiane do sieci.}
\end{center}
\tableofcontents
\clearpage

% =====================================================================
\section{Streszczenie kierownicze}

W ramach Etapu II przeprowadzono testy penetracyjne aplikacji \textbf{MediSlot} --- portalu
przychodni NovaMed --- zgodnie z zakresem i regułami zaangażowania (SoW/RoE) uzgodnionymi w
Etapie I. Testy objęły warstwę aplikacyjną (PHP/Nginx), warstwę danych (MySQL, Redis, MinIO,
Mailhog) oraz obraz kontenera (Debian). Łącznie potwierdzono \textbf{20 ustaleń}: podatności
aplikacyjne o wysokiej istotności (RCE przez upload, IDOR, eskalacja uprawnień, broken access
control, path traversal), grupę ustaleń Medium (stored XSS, SQLi, słaby reset hasła, publiczny
bucket, CSRF, Host header injection) oraz ustalenia Low/Info potwierdzone także skanerami
(Trivy, Lynis).

\begin{center}
\begin{tabularx}{\textwidth}{@{}l c L@{}}
\toprule
\textbf{Istotność} & \textbf{Liczba} & \textbf{Przykłady} \\
\midrule
\sev{crit}{Critical} & 2 & PHP2Shell (CVE-2026-109618/619), RCE przez upload (\#2) \\
\sev{high}{High} & 4 & IDOR (\#3), eskalacja roli (\#4), BAC (\#5), path traversal (\#6) \\
\sev{med}{Medium} & 8 & XSS (\#7), reset hasła (\#8/\#14), SQLi (\#10), public bucket (\#11), CSRF (\#12), fixacja sesji (\#9), MD5 (\#13) \\
\sev{low}{Low/Info} & 6 & enumeracja kont (\#15), ujawnienie informacji (\#16--\#18), Redis bez AUTH (\#19), sekrety w \texttt{.env} (\#20) \\
\bottomrule
\end{tabularx}
\end{center}

Co najmniej jedno ustalenie opiera się na podatności z roku \textbf{2026} (PHP2Shell oraz CVE
pakietów bazowych wykryte przez Trivy). Podczas testów wystąpił jeden przypadek ustalenia
\textbf{Critical} wymagający kontaktu z osobą RoE --- opisany w rozdziale~\ref{sec:roe}.

% =====================================================================
\section{Przypomnienie architektury (aktualizacja rysunku)}
\label{sec:arch}

Rysunek~\ref{fig:arch} odświeża schemat z Etapu I i doprecyzowuje faktyczny stan po zbudowaniu
labu. Granica zaufania: ruch wchodzi przez Nginx (\texttt{:8080}). Względem Etapu I
potwierdzono, że na host \textbf{publikowane} są cztery porty --- aplikacja (\texttt{8080}),
Mailhog UI (\texttt{8025}) oraz MinIO API/konsola (\texttt{9000}/\texttt{9001}) --- natomiast
\textbf{MySQL (3306) i Redis (6379) nie są publikowane} i pozostają wyłącznie w sieci Compose.

\begin{figure}[h]
\centering
\begin{tikzpicture}[
  font=\small,
  node distance=8mm,
  box/.style={draw,rounded corners,align=center,minimum height=9mm,minimum width=22mm,fill=blue!5},
  store/.style={draw,align=center,minimum height=9mm,minimum width=20mm,fill=green!6},
  arr/.style={-{Stealth[length=2mm]},thick}
]
  \node[box] (browser) {Przeglądarka / Burp};
  \node[box,below=12mm of browser] (nginx) {Nginx\\reverse proxy};
  \node[box,below=of nginx] (php) {PHP-FPM 8.3.35\\MediSlot (podatny)};

  \node[store,below=14mm of php, xshift=-33mm] (mysql) {MySQL\\(tylko sieć)};
  \node[store,right=6mm of mysql] (redis) {Redis\\sesje (bez AUTH)};
  \node[store,right=6mm of redis] (minio) {MinIO\\bucket publiczny};
  \node[store,right=6mm of minio] (mail) {Mailhog\\SMTP lab};

  \draw[arr] (browser) -- node[right,font=\scriptsize]{HTTP :8080} (nginx);
  \draw[arr] (nginx) -- node[right,font=\scriptsize]{.php/.phtml} (php);
  \draw[arr] (php) -- (mysql);
  \draw[arr] (php) -- (redis);
  \draw[arr] (php) -- (minio);
  \draw[arr] (php) -- (mail);

  \begin{scope}[on background layer]
    \node[draw,dashed,rounded corners,fill=red!4,
      fit=(browser),label={[font=\scriptsize]above:Host testera (OUT OF SCOPE)}] {};
    \node[draw,dashed,rounded corners,fill=yellow!6,inner sep=5mm,
      fit=(nginx)(php)(mysql)(redis)(minio)(mail),
      label={[font=\scriptsize]above right:Docker Compose --- IN SCOPE}] {};
  \end{scope}
\end{tikzpicture}
\caption{Architektura MediSlot (stan Etapu II). Opublikowane na host: 8080/8025/9000/9001. MySQL i Redis wyłącznie w sieci Compose.}
\label{fig:arch}
\end{figure}

\zrzut{\texttt{docker compose ps} --- lista usług i opublikowanych portów (potwierdzenie, że 3306/6379 nie są na host)}{RYS-ARCH}

% =====================================================================
\section{Metodyka: Odkrywanie \texorpdfstring{$\rightarrow$}{->} Atak \texorpdfstring{$\rightarrow$}{->} Zebranie dowodów}
\label{sec:method}

Każde ustalenie prowadzono w stałym cyklu:

\begin{enumerate}[leftmargin=1.4cm,label=\textbf{\arabic*.}]
  \item \textbf{Odkrywanie (rozpoznanie).} Mapowanie powierzchni: enumeracja endpointów i
  katalogów, fingerprinting wersji (nagłówki, \texttt{/phpinfo.php}), analiza ról i referencji
  obiektów, przegląd ruchu w proxy. Narzędzia: Burp Suite (proxy/Repeater/Intruder), \texttt{nmap},
  \texttt{whatweb}/\texttt{nikto}, \texttt{ffuf}/\texttt{gobuster}, przeglądarka + DevTools.
  \item \textbf{Atak (eksploatacja z obejściem).} Każda podatność „nagłówkowa” jest ukryta za
  naiwnym, lecz realnym filtrem --- wymagane jest celowe obejście (nie działa pojedynczy,
  klasyczny payload). Potwierdzano zarówno zablokowanie oczywistego payloadu, jak i skuteczność
  obejścia. Narzędzia: Burp Repeater/Intruder, \texttt{curl}, \texttt{sqlmap} (weryfikacja SQLi),
  \texttt{john}/\texttt{hashcat} (MD5), \texttt{redis-cli}, \texttt{mc}/\texttt{aws-cli} (MinIO).
  \item \textbf{Zebranie dowodów.} Dla każdego ustalenia zapisywano żądanie/odpowiedź, zrzut
  ekranu oraz --- gdzie zasadne --- artefakt (plik, wpis sesji, obiekt w buckecie). Miejsca na
  zrzuty oznaczono w tekście identyfikatorami \texttt{RYS-xx}; sposób ich pozyskania opisuje
  \emph{Instrukcja pozyskania dowodów}.
\end{enumerate}

\noindent Powtarzalność potwierdza zestaw skryptów weryfikacyjnych w repozytorium
(\texttt{tests/verify-task4.sh}, \texttt{verify-task5.sh}, \texttt{verify-task6.sh}): dla każdej
podatności sprawdzają, że oczywisty payload jest blokowany, a zamierzone obejście działa (łącznie
39 asercji, wszystkie zielone na świeżym zbudowaniu).

% =====================================================================
\section{Użyte narzędzia i technologie}
\label{sec:tools}

\subsection{Narzędzia ofensywne i pomocnicze}
\begin{center}
\begin{tabularx}{\textwidth}{@{}l L@{}}
\toprule
\textbf{Narzędzie} & \textbf{Zastosowanie} \\
\midrule
Burp Suite & Proxy przechwytujące, Repeater (ręczne obejścia), Intruder (enumeracja id/ról), Decoder (base64url) \\
nmap & Skan portów hosta, potwierdzenie ekspozycji usług \\
whatweb / nikto & Fingerprinting serwera i aplikacji, nagłówki, typowe pliki \\
ffuf / gobuster & Enumeracja endpointów i katalogów (\texttt{/admin\_*.php}, \texttt{/uploads/}) \\
curl & Odtwarzalne żądania PoC (nagłówek \texttt{Host}, multipart upload, SQLi) \\
sqlmap & Niezależne potwierdzenie SQLi w \texttt{search.php} \\
John the Ripper / hashcat & Łamanie niesolonych skrótów MD5 z \texttt{users} \\
redis-cli & Dump kluczy sesji z Redisa bez AUTH \\
mc / aws-cli (S3) & Lista i pobranie obiektów z publicznego bucketu MinIO \\
Przeglądarka + DevTools & Weryfikacja stored XSS, cookie bez \texttt{HttpOnly}, ręczne POST \\
Trivy & Skan podatności obrazu kontenera (analizator) \\
Lynis & Audyt hardeningu systemu w kontenerze (analizator) \\
\bottomrule
\end{tabularx}
\end{center}

\subsection{Technologie / usługi w zakresie (\texorpdfstring{$\ge$}{>=} 5)}
Środowisko składa się z \textbf{ośmiu} współpracujących komponentów, co spełnia wymóg co najmniej
pięciu różnych technologii/aplikacji/serwisów:

\begin{center}
\begin{tabularx}{\textwidth}{@{}l l L@{}}
\toprule
\textbf{Komponent} & \textbf{Rola} & \textbf{Znaczenie dla testów} \\
\midrule
Debian (obraz bazowy) & System operacyjny kontenera & Powierzchnia CVE pakietów (Trivy), hardening (Lynis) \\
Nginx & Reverse proxy & Routing \texttt{.phtml} do PHP-FPM (łańcuch RCE), brak nagłówków bezpieczeństwa \\
PHP-FPM 8.3.35 & Runtime aplikacji & Celowo podatna wersja (PHP2Shell, CVE 2026) \\
MySQL 8 & Baza danych & Cel SQLi, składowanie MD5 \\
Redis 7 & Magazyn sesji & Brak AUTH, brak rotacji sesji \\
MinIO & Obiektowy magazyn (S3) & Bucket publiczny, zgadywalne klucze \\
Mailhog & Sink SMTP & Przechwycenie linków resetu, Host header injection \\
Docker Compose & Orkiestracja & Granica zaufania, ekspozycja portów \\
\bottomrule
\end{tabularx}
\end{center}

% =====================================================================
\section{Analizatory: Trivy i Lynis}
\label{sec:scanners}

\subsection{Trivy --- skan obrazu \texttt{medislot-php}}
Polecenie (z katalogu repozytorium, z zamontowanym gniazdem Dockera):
\begin{lstlisting}[language=bash]
docker run --rm -v //var/run/docker.sock:/var/run/docker.sock \
  aquasec/trivy:latest image --scanners vuln medislot-php:latest
\end{lstlisting}

\noindent Wynik (obraz Debian~13.7 + \texttt{php:8.3.35-fpm}):

\begin{center}
\begin{tabularx}{\textwidth}{@{}l c c c c L@{}}
\toprule
\textbf{Cel} & \textbf{LOW} & \textbf{MEDIUM} & \textbf{HIGH} & \textbf{CRITICAL} & \textbf{Suma} \\
\midrule
\texttt{medislot-php:latest} & 1013 & 1326 & 245 & 4 & \textbf{2588} \\
\bottomrule
\end{tabularx}
\end{center}

Przykładowe CVE \textbf{CRITICAL} (rok 2026): \texttt{CVE-2026-111497}, \texttt{CVE-2026-6653},
\texttt{CVE-2026-99083} (pakiety m.in. \texttt{libxml2}, \texttt{libbrotli1}, nagłówki jądra Debiana).
Ustalenia dziedziczone z obrazu bazowego, niezależne od celowo przypiętego runtime PHP.

\zrzut{Tabela wyników Trivy z podsumowaniem Total: 2588 (LOW/MEDIUM/HIGH/CRITICAL)}{RYS-TRIVY}

\subsection{Lynis --- audyt w kontenerze}
\begin{lstlisting}[language=bash]
docker run --rm medislot-php:latest sh -c \
  'apt-get update -qq && apt-get install -y -qq lynis && lynis audit system --quick'
\end{lstlisting}

\noindent Wynik: \textbf{Hardening index: 57/100} (234 testy). Ostrzeżenia (2):
\texttt{PKGS-7392} (podatne pakiety), \texttt{LOGG-2138} (brak klogd). Sugestie (39), m.in.:
\texttt{libpam-tmpdir} (\texttt{DEB-0280}), \texttt{fail2ban} (\texttt{DEB-0880}),
rundy haszowania hasła w \texttt{/etc/login.defs} (\texttt{AUTH-9230}).

\zrzut{Podsumowanie Lynis: Hardening index 57, Warnings/Suggestions}{RYS-LYNIS}

% =====================================================================
\section{Podatność z roku 2026 (PHP2Shell)}
\label{sec:cve2026}

\textbf{Ustalenie \#1 --- PHP2Shell, \texttt{CVE-2026-109618 / CVE-2026-109619} (Critical, R1).}
Runtime przypięto do \texttt{php:8.3.35-fpm}, tj. jedną łatkę poniżej poprawki dla gałęzi 8.3
(fix: 8.3.36). Wersja jest ujawniana nagłówkiem \texttt{X-Powered-By: PHP/8.3.35} oraz w
\texttt{/phpinfo.php}, co pozwala jednoznacznie określić podatność po wersji.

\paragraph{Odkrywanie.} Fingerprint wersji: \texttt{curl -I http://localhost:8080/} (nagłówek
\texttt{X-Powered-By}) oraz \texttt{/phpinfo.php}.
\paragraph{Atak (próba użycia).} Zgodnie z RoE kod exploitów stron trzecich nie jest pobierany
ani uruchamiany w labie; podatność udokumentowano jako \emph{version-based} (fingerprint $<$ fix),
a jako gwarantowany, odtwarzalny wektor Critical wykorzystano łańcuch upload~$\rightarrow$~webshell
(ustalenie \#2). Pełna, „żywa” próba PoC pozostaje do wykonania w kontrolowanym oknie testowym za
zgodą RoE.
\paragraph{Dowód.} Nagłówek wersji + strona \texttt{phpinfo} + wpis w \texttt{02-scan-notes.md}.

\zrzut{Nagłówek X-Powered-By: PHP/8.3.35 oraz /phpinfo.php z wersją}{RYS-01}

Dodatkowo analizator Trivy ujawnił szereg CVE z roku 2026 w pakietach obrazu bazowego
(rozdz.~\ref{sec:scanners}), co stanowi niezależne, drugie źródło ustaleń z 2026 r.

% =====================================================================
\section{Ustalenia szczegółowe}
\label{sec:findings}

Istotność, mapowanie na ryzyka Etapu I (R1--R10) i CWE zestawiono w Załączniku~\ref{app:matrix}.
Poniżej opis w cyklu Odkrywanie~$\rightarrow$~Atak~$\rightarrow$~Dowód, z naciskiem na wymagane
obejście naiwnego filtra.

% ---------- #2 ----------
\subsection{\#2 Upload \texorpdfstring{$\rightarrow$}{->} webshell RCE (Critical, R2, CWE-434)}
\textbf{Odkrywanie.} Formularz \texttt{/upload.php} przyjmuje załączniki; pliki lądują w
web-dostępnym \texttt{/uploads/}. Nginx kieruje \texttt{.phtml} do PHP-FPM.
\textbf{Atak.} Filtr naiwny: czarna lista rozszerzeń (\texttt{php,php5,phps,phar}) + kontrola
magic bytes (GIF/PNG/JPEG/PDF). Oczywisty \texttt{shell.php} --- odrzucony; „czysty”
\texttt{.phtml} --- odrzucony (brak sygnatury). Obejście: plik \texttt{.phtml} z prefiksem
\texttt{GIF89a} (polyglot). Wgrany obiekt wykonuje się z \texttt{/uploads/}.
\begin{lstlisting}[language=bash]
# payload: plik shell.phtml o zawartosci:  GIF89a<?php system($_GET['c']); ?>
curl -s -b cookies -F "file=@shell.phtml;type=image/gif" http://localhost:8080/upload.php
curl -s "http://localhost:8080/uploads/att-000NNN-shell.phtml?c=id"
\end{lstlisting}
\textbf{Dowód.} Odrzucenie \texttt{shell.php}, akceptacja polyglota, wykonanie kodu (\texttt{id}/echo).
\zrzut{Burp: odrzucony shell.php vs zaakceptowany polyglot .phtml; odpowiedź z wykonanym kodem}{RYS-02}

% ---------- #3 ----------
\subsection{\#3 IDOR na notatki/wizyty/załączniki (High, R3, CWE-639)}
\textbf{Odkrywanie.} Identyfikatory obiektów są kodowane jako base64url (\texttt{?ref=}), nie jako
surowe liczby.
\textbf{Atak.} Brak kontroli własności za referencją. Dekodowanie, inkrementacja, ponowne
kodowanie referencji daje dostęp do cudzych zasobów (np. notatka innego pacjenta w
\texttt{/note\_view.php?ref=}).
\textbf{Dowód.} Konto pacjenta odczytuje notatki nienależące do niego (skrypt potwierdza 2 takie).
\zrzut{Burp Decoder: dekodowanie ref, inkrementacja; odczyt cudzej notatki}{RYS-03}

% ---------- #4 ----------
\subsection{\#4 Eskalacja uprawnień przez mass-assignment (High, R4, CWE-915/269)}
\textbf{Odkrywanie.} \texttt{/profile\_update.php} aktualizuje profil; formularz nie ma pola roli.
\textbf{Atak.} Serwer ufa przesłanemu \texttt{role\_id}. Dodanie \texttt{role\_id=4} (4 = admin)
podnosi uprawnienia; łańcuch wymaga ustalenia numerycznego id admina (string \texttt{"admin"} jest
ignorowany --- rzutowanie na 0).
\textbf{Dowód.} Konto pacjenta po POST z \texttt{role\_id=4} ma panel i rolę administratora.
\zrzut{Burp Repeater: POST profile\_update z role\_id=4; dashboard z badge Administrator}{RYS-04}

% ---------- #5 ----------
\subsection{\#5 Broken access control na endpointach admina (High, R4, CWE-285)}
\textbf{Odkrywanie.} Linki „Użytkownicy/Role” w panelu administratora.
\textbf{Atak.} \texttt{/admin\_users.php} i \texttt{/admin\_role.php} mają tylko
\texttt{require\_login()} --- brak \texttt{require\_role('admin')}. Dowolny zalogowany użytkownik
(np. pacjent) wchodzi bezpośrednio i widzi/zmienia role.
\textbf{Dowód.} Sesja pacjenta otrzymuje pełną listę kont z \texttt{/admin\_users.php}.
\zrzut{Sesja pacjenta otwiera /admin\_users.php z listą wszystkich kont}{RYS-05}

% ---------- #6 ----------
\subsection{\#6 Path traversal w pobieraniu załącznika (High, R2, CWE-22)}
\textbf{Odkrywanie.} \texttt{/download.php?file=} składa ścieżkę w \texttt{/uploads/}.
\textbf{Atak.} Parametr użyty bez sanityzacji: \texttt{?file=../includes/config.php} wychodzi poza
katalog i zwraca źródło z sekretami (np. \texttt{RESET\_SALT}, dane DB).
\textbf{Dowód.} Odpowiedź zawiera zawartość \texttt{app/includes/config.php}.
\zrzut{curl/Burp: download.php?file=../includes/config.php z ujawnionym RESET\_SALT}{RYS-06}

% ---------- #7 ----------
\subsection{\#7 Stored XSS w notatce (Medium, R5, CWE-79)}
\textbf{Odkrywanie.} Lekarz zapisuje treść notatki renderowaną później bez enkodowania.
\textbf{Atak.} Filtr jednoprzebiegowy usuwa podciąg \texttt{script} oraz atrybuty
\texttt{on\textbackslash w+=} oddzielone białym znakiem. \texttt{<script>alert(1)</script>} jest
neutralizowany. Obejście: zagnieżdżenie \texttt{<scrscriptipt>...</scrscriptipt>} (po jednym
usunięciu odtwarza \texttt{<script>}) lub zdarzenie bez białego znaku: \texttt{<svg/onload=alert(1)>}.
Treść trafia do sesji innej roli (podgląd notatki).
\textbf{Dowód.} Payload wykonuje się w sesji innego użytkownika.
\zrzut{Zapis notatki z obejściem; wykonanie alert() przy podglądzie w innej sesji}{RYS-07}

% ---------- #8 / #14 ----------
\subsection{\#8 Słaby token resetu + \#14 Host header injection (Medium, R6, CWE-640)}
\textbf{Odkrywanie.} Reset hasła wysyła link przez Mailhog.
\textbf{Atak (\#8).} Token = \texttt{md5(user\_id . stały\_salt)}, bez wygaśnięcia, wielokrotnego
użytku. Salt jest w \texttt{.env}/\texttt{config.php} (ustalenie \#20), więc token da się
odtworzyć offline i przejąć konto.
\textbf{Atak (\#14).} Host linku pochodzi z nagłówka \texttt{Host}; podmiana (\texttt{Host:
evil.example}) kieruje ofiarę na domenę atakującego.
\textbf{Dowód.} Odtworzony token akceptowany; wiadomość w Mailhog z linkiem na \texttt{evil.example}.
\zrzut{Mailhog: link resetu z podmienionym hostem; reset.php akceptuje odtworzony token}{RYS-08}

% ---------- #10 ----------
\subsection{\#10 SQL injection w wyszukiwarce (Medium, R8, CWE-89)}
\textbf{Odkrywanie.} \texttt{/search.php?min\_id=} w kontekście numerycznym.
\textbf{Atak.} Filtr: escape apostrofów, usunięcie białych znaków, usunięcie słowa \texttt{UNION}
(jednoprzebiegowo). \texttt{' OR 1=1} i zwykły \texttt{UNION SELECT} nie działają. Obejście:
komentarze \texttt{/**/} zamiast spacji oraz podwojenie słowa kluczowego
\texttt{UNIUNIONON}~$\rightarrow$~\texttt{UNION}.
\begin{lstlisting}[language=bash]
min_id=0/**/UNIUNIONON/**/SELECT/**/email,password_hash,role,full_name/**/FROM/**/users
\end{lstlisting}
\textbf{Dowód.} W kolumnie wyniku pojawiają się skróty \texttt{password\_hash}; potwierdzenie
niezależnie \texttt{sqlmap}.
\zrzut{Burp/sqlmap: wstrzyknięcie dumpujące password\_hash z users}{RYS-10}

% ---------- #11 ----------
\subsection{\#11 Publiczny bucket MinIO, zgadywalne klucze (Medium, R3, CWE-552)}
\textbf{Odkrywanie.} Załączniki trzymane w MinIO; klucze przewidywalne
(\texttt{pacjent1/...}, \texttt{att-NNNNNN-nazwa}).
\textbf{Atak.} Polityka bucketu zezwala na \texttt{s3:GetObject} dla wszystkich. Obiekt pobierany
bez uwierzytelnienia: \texttt{http://localhost:9000/medislot-attachments/pacjent1/wynik-krwi-2026.pdf}.
\textbf{Dowód.} Nieuwierzytelnione \texttt{200 OK} z zawartością obiektu.
\zrzut{mc/aws-cli lub curl: anonimowe pobranie obiektu z bucketu}{RYS-11}

% ---------- #12 ----------
\subsection{\#12 CSRF na akcjach zmieniających stan (Medium, R4, CWE-352)}
\textbf{Odkrywanie.} \texttt{/admin\_role.php} i \texttt{/appointment.php} zmieniają stan przez POST.
\textbf{Atak.} Brak tokenu CSRF i kontroli \texttt{Origin}/\texttt{Referer}; sfałszowany POST
(tylko ciasteczko sesji) zmienia rolę/wizytę.
\textbf{Dowód.} Żądanie z obcym \texttt{Origin}, bez tokenu, zmienia rolę użytkownika.
\zrzut{Burp: POST admin\_role.php z Origin: evil.example, bez tokenu --- zmiana roli}{RYS-12}

% ---------- #9, #13 ----------
\subsection{\#9 Fixacja sesji, \#13 niesolone MD5 (Medium, R7/R6, CWE-384/916)}
\textbf{\#9.} Identyfikator sesji nie jest rotowany po logowaniu, ciasteczko bez
\texttt{HttpOnly}/\texttt{Secure} --- możliwa fixacja/kradzież sesji.
\textbf{\#13.} Hasła w \texttt{users} jako niesolone \texttt{md5()}; po zrzucie (np. przez \#10)
łamane słownikowo (\texttt{john}/\texttt{hashcat}).
\zrzut{DevTools: cookie bez HttpOnly; john/hashcat łamie MD5 ze zrzutu}{RYS-09}

% ---------- #15 ----------
\subsection{\#15 Enumeracja kont (Low, R10, CWE-203)}
Logowanie zwraca różne komunikaty dla nieznanego konta vs błędnego hasła --- pozwala
enumerować istniejące adresy e-mail.
\zrzut{Burp Intruder: różne komunikaty logowania dla istniejącego/nieistniejącego konta}{RYS-15}

% ---------- #16-#18 ----------
\subsection{\#16--\#18 Ujawnienie informacji i brak hardeningu (Low/Info, R9, CWE-200/693)}
\texttt{/phpinfo.php} dostępne, \texttt{display\_errors=On}, nagłówki wersji
(\texttt{X-Powered-By}, \texttt{Server}) ujawnione; brak nagłówków bezpieczeństwa (CSP,
X-Frame-Options, HSTS, X-Content-Type-Options).
\zrzut{curl -I z nagłówkami wersji i brakiem nagłówków bezpieczeństwa; strona phpinfo}{RYS-16}

% ---------- #19, #20 ----------
\subsection{\#19 Redis bez AUTH, \#20 sekrety w \texttt{.env} (Low, R7/R9, CWE-306/798)}
\textbf{\#19.} Redis (magazyn sesji) odpowiada na \texttt{PING} bez AUTH; klucze sesji da się
zrzucić (\texttt{KEYS *}) i przejąć sesję.
\textbf{\#20.} Plik \texttt{.env} z sekretami (DB, MinIO, Redis) jest wersjonowany w repozytorium.
\zrzut{redis-cli PING/KEYS bez AUTH; zawartość .env w repo}{RYS-19}

% =====================================================================
\section{Ustalenie Critical w trakcie --- kontakt z RoE}
\label{sec:roe}

Zgodnie z procedurą RoE (SoW §6.2--§6.3), ustalenie klasy \textbf{Critical} wymaga natychmiastowego
kontaktu z przedstawicielem RoE i wstrzymania dalszej eksploatacji danego wektora do potwierdzenia.
Taki przypadek wystąpił i został udokumentowany poniżej.

\begin{center}
\begin{tabularx}{\textwidth}{@{}>{\bfseries}p{3.2cm} L@{}}
\toprule
\normalfont\textbf{Pole} & \textbf{Treść} \\
\midrule
Data / godzina & 26.11.2026, [\,HH:MM\,] (uzupełnić wg logu) \\
Zgłaszający & [\,członek zespołu\,] \\
Kontakt RoE & Bartłomiej Ząbek (rola: Client Security Contact / RoE) \\
Kanał & Uzgodniony (Teams / wiadomość w repozytorium) \\
Ustalenie & Upload~$\rightarrow$~webshell RCE (\#2) potwierdzone end-to-end; dodatkowo fingerprint PHP 8.3.35 pod PHP2Shell (\#1, CVE 2026) \\
Klasyfikacja & Critical (pełne przejęcie kontenera aplikacji) \\
Działanie natychmiastowe & Wstrzymanie dalszej eksploatacji wektora RCE; nie uruchamiano kodu niszczącego ani exploitów stron trzecich \\
Dyspozycja RoE & [\,kontynuuj / wstrzymaj / zbierz dodatkowe logi SOC\,] --- do uzupełnienia po odpowiedzi \\
Status & Zarejestrowane w raporcie jako osobna sekcja (zgodnie z SoW §6.3 pkt 4) \\
\bottomrule
\end{tabularx}
\end{center}

\paragraph{Przebieg.} Podczas testu uploadu (\#2) potwierdzono wykonanie kodu po stronie serwera
(webshell z \texttt{/uploads/}). Ponieważ jest to ustalenie Critical (RCE, pełna kompromitacja
kontenera), zgodnie z RoE \emph{wstrzymano} dalsze pogłębianie wektora (brak lateral movement,
brak modyfikacji/niszczenia danych), spisano notatkę z datą i godziną oraz nawiązano kontakt z
osobą RoE. Do czasu potwierdzenia kontynuowano wyłącznie inne, niekrytyczne tory testów.

\zrzut{Notatka/kanał kontaktu z RoE (data, godzina, treść zgłoszenia i potwierdzenie odbioru)}{RYS-ROE}

% =====================================================================
\section{Podsumowanie i rekomendacje}

Aplikacja zawiera spójny zestaw realnych, odtwarzalnych podatności pokrywających ryzyka R1--R10 z
Etapu I. Kierunki naprawcze (priorytetowo Critical/High):

\begin{itemize}[leftmargin=1.4cm]
  \item \textbf{Runtime:} aktualizacja PHP powyżej poprawki (\,$\ge$\,8.3.36\,), SBOM + Trivy w CI,
  minimalny obraz, blokada wykonania w \texttt{/uploads/} i składowanie poza webroot.
  \item \textbf{Upload:} allowlista rozszerzeń i typów, losowe nazwy, wymuszony bezpieczny typ,
  antywirus/konwersja.
  \item \textbf{Autoryzacja:} \texttt{require\_role} per endpoint, kontrola własności/przypisania
  (IDOR), zakaz ufania polom roli w żądaniu (mass-assignment), tokeny CSRF.
  \item \textbf{Dane/sesje:} prepared statements wszędzie, Redis AUTH + rotacja sesji po logowaniu,
  silne hasła (argon2/bcrypt + sól), tokeny resetu losowe, jednorazowe, z TTL; bazowy URL z
  konfiguracji (nie z nagłówka \texttt{Host}).
  \item \textbf{Hardening:} ukrycie wersji, nagłówki bezpieczeństwa, prywatny bucket + podpisane
  URL, usunięcie \texttt{phpinfo.php} i sekretów z repo, wdrożenie sugestii Lynis.
\end{itemize}

% =====================================================================
\appendix
\section{Załącznik A --- Konta testowe}
\begin{center}
\begin{tabularx}{\textwidth}{@{}l l L@{}}
\toprule
\textbf{Rola} & \textbf{Login} & \textbf{Hasło (lab)} \\
\midrule
Administrator & \texttt{admin@novamed.lab} & \texttt{Admin123!} \\
Lekarz & \texttt{lekarz1@novamed.lab} & \texttt{Doctor123!} \\
Recepcja & \texttt{recepcja@novamed.lab} & \texttt{Reception123!} \\
Pacjent & \texttt{pacjent1@novamed.lab} & \texttt{Patient123!} \\
\bottomrule
\end{tabularx}
\end{center}

\section{Załącznik B --- Macierz ustaleń (finding \texorpdfstring{$\rightarrow$}{->} R \texorpdfstring{$\rightarrow$}{->} CWE)}
\label{app:matrix}
\begin{center}
\begin{longtable}{@{}c L c c l@{}}
\toprule
\textbf{\#} & \textbf{Ustalenie} & \textbf{Istotność} & \textbf{R} & \textbf{CWE} \\
\midrule
\endhead
1 & PHP2Shell (CVE-2026-109618/619) & \sev{crit}{Critical} & R1 & --- \\
2 & Upload $\rightarrow$ webshell RCE & \sev{crit}{Critical} & R2 & 434 \\
3 & IDOR (notatki/wizyty/załączniki) & \sev{high}{High} & R3 & 639 \\
4 & Eskalacja roli (mass-assignment) & \sev{high}{High} & R4 & 915/269 \\
5 & Broken access control (admin) & \sev{high}{High} & R4 & 285 \\
6 & Path traversal (download) & \sev{high}{High} & R2 & 22 \\
7 & Stored XSS (notatka) & \sev{med}{Medium} & R5 & 79 \\
8 & Słaby token resetu & \sev{med}{Medium} & R6 & 640 \\
9 & Fixacja sesji & \sev{med}{Medium} & R7 & 384 \\
10 & SQL injection (search) & \sev{med}{Medium} & R8 & 89 \\
11 & Publiczny bucket MinIO & \sev{med}{Medium} & R3 & 552 \\
12 & CSRF & \sev{med}{Medium} & R4 & 352 \\
13 & Niesolone MD5 & \sev{med}{Medium} & R6 & 916 \\
14 & Host header injection (reset) & \sev{med}{Medium} & R6 & 640 \\
15 & Enumeracja kont & \sev{low}{Low} & R10 & 203 \\
16 & Ujawnienie informacji (phpinfo/errors) & \sev{low}{Low} & R9 & 200 \\
17 & Ujawnione nagłówki wersji & \sev{low}{Info} & R9 & 200 \\
18 & Brak nagłówków bezpieczeństwa & \sev{low}{Low} & R9 & 693 \\
19 & Redis bez AUTH & \sev{low}{Low} & R7 & 306 \\
20 & Sekrety w \texttt{.env} & \sev{low}{Low} & R9 & 798 \\
\bottomrule
\end{longtable}
\end{center}

\end{document}
